\Page{5}{\textenglish{01 / DESIGN}}{הקולקציה שעוד לא רואים}{42}{50}

\Lead{קולקציה מגיעה למסלול כשהיא כבר יודעת מה היא רוצה להיות. בסטודיו, כמה חודשים קודם, ההחלטות מסודרות פחות. לגוון שהמעצבת אוהבת יש שתי גרסאות קרובות. בתוך אותו קו יש צללית אחת שאפשר לדמיין מיד ברחוב, ואחרת שעשויה לשנות את האופי של הקולקציה כולה. כדי לדעת מה שווה לפתח, הצוות צריך לתת לכל רעיון צורה מספקת כדי שאפשר יהיה לשפוט אותו.}

הדרך המקובלת מתחילה בהשראה ובסקיצות. השרטוט הטכני מבהיר את הפרטים שהצוות צריך לפתח, הגזרה נותנת להם מידות וצורה, ובמדידה של הדוגמה רואים כיצד הכול עובד יחד על הגוף. לקראת הייצור מרכזים בתיק המוצר את המידות, החומרים וההוראות שהמפעל צריך. ברגע שהצוות מתחייב לדוגמה, הוא משקיע באפשרות אחת זמן וחומר. ההשקעה הזאת מסבירה למה גם רעיון מעניין יכול להישאר בשלב הסקיצה. \SourceRef{1}

\vspace{5mm}\begin{english}\centering
\begin{tikzpicture}[
  stage/.style={draw=Silver,fill=PaperSteel,minimum width=47mm,minimum height=17mm,
    align=center,font=\fontsize{12}{16}\selectfont,inner sep=2.5mm},
  flow/.style={draw=Warm,line width=.65pt,-{Latex[length=2mm]}}]
  \node[stage] (sketch) at (11.6,0) {\texthebrew{רעיון ושרטוט}};
  \node[stage] (pattern) at (5.8,0) {\texthebrew{גזרה ודוגמה}};
  \node[stage] (fitting) at (0,0) {\texthebrew{מדידה ומפרט}};
  \draw[flow] (sketch.west) -- (pattern.east);
  \draw[flow] (pattern.west) -- (fitting.east);
\end{tikzpicture}\par
\end{english}
\Caption{שלוש תחנות בתהליך פיתוח הבגד, מן הרעיון ועד למפרט שהמפעל מקבל.}



{\fontsize{13.8}{21.3}\selectfont

\StoryColumns{במדידה מתברר אילו החלטות צריכות להשתנות. כתף שנראתה משוחררת בשרטוט יכולה למשוך את החלק העליון כשהיד עולה. נפח שיפה בעמידה יכול להכביד בישיבה. התיקון עשוי לשנות את הגזרה, את המפרט או את בחירת הבד. בסטודיו קטן אותה מעצבת יכולה ללוות את כל העבודה. בצוות גדול מנהלת פיתוח ותדמיתנית עשויות לקחת אחריות על חלקים שונים שלה. לכל אחת כדאי לקבל גרסה שמבהירה מה הוסכם ומה עדיין פתוח. \SourceRef{1}}{כאן נוצרת שאלת הקולקציה. אם כבר השקענו זמן בדוגמה שכולם מבינים, קל יותר לשפר אותה מאשר לחזור לרעיון שקשה להסביר. היכולת להעמיד שתי אפשרויות זו מול זו מוקדם יותר יכולה לתת מקום גם לפריט הפחות מובן מאליו.}

\par}
